% Compile with pdfLaTeX (no external files or bibliography processor needed).
% Keep the submitted proposal to ONE page, including references.
\documentclass[11pt, letterpaper]{article}
\usepackage[T1]{fontenc}
\usepackage[utf8]{inputenc}
\usepackage[margin=0.75in, paperheight=11in]{geometry}
\usepackage{newtxtext,newtxmath}
\usepackage{microtype}
\usepackage{xcolor}
\usepackage{titlesec}
\usepackage{hyperref}

\definecolor{ProposalBlue}{HTML}{40619F}
\definecolor{SecondaryText}{HTML}{555555}
\hypersetup{colorlinks=true,urlcolor=ProposalBlue,linkcolor=ProposalBlue,
  pdfauthor={},pdftitle={New-Idea Proposal for a Paper Presentation}}
\pagestyle{empty}
\setlength{\parindent}{0pt}
\setlength{\parskip}{3pt}
\setlength{\emergencystretch}{1em}
\titleformat{\section}{\normalfont\bfseries\color{ProposalBlue}}{}{0pt}{}
\titlespacing*{\section}{0pt}{7pt}{2pt}
\newcommand{\proposalheader}[4]{%
  {\small\sffamily\bfseries\color{ProposalBlue}
   PAPER PRESENTATION\enspace\textbullet\enspace NEW-IDEA PROPOSAL}\par
  \vspace{3pt}
  {\fontsize{16}{18}\selectfont\bfseries #1\par}
  \vspace{4pt}
  \textbf{Student(s):} #2\hfill\textbf{Presentation:} #3\par
  \textbf{Based on:} #4\par
  \vspace{3pt}
  {\color{ProposalBlue}\hrule height 0.6pt}
  \vspace{4pt}
}
\newcommand{\guide}[1]{{\color{SecondaryText}#1}}
\newcommand{\wordbudget}[1]{{\normalfont\small\color{SecondaryText}\hfill #1 words}}
\newcommand{\referencesheading}{%
  \par\vspace{4pt}
  {\color{ProposalBlue}\hrule height 0.4pt}
  \vspace{3pt}
}

\begin{document}

\proposalheader{[Specific title of your proposed research idea]}
  {[Your name(s)], Dimitri Lyon}{10/28/2026}
  {[Paper title, authors, and year]}

\guide{\textit{Propose one testable extension of the presented paper, not just a summary.
Replace the bracketed prompts. Keep one page, including references, with 11-point text
and the existing margins. Suggested length: 420--520 words.}}

\section*{1. Background and motivation\wordbudget{50--60}}
[Explain the problem the paper addresses, why it matters, and the main contribution
that your idea builds on. Give enough context for a classmate who has not read the
paper. Cite the relevant method or result; do not summarize every section.]

\section*{2. Gap and research question\wordbudget{50--60}}
[Identify one limitation, missing comparison, or untested assumption. Point to a
specific section, figure, experiment, or equation when possible. Distinguish the
paper's stated limitation from your own observation. End with one focused question:
``Can [specific change] improve [outcome] under [conditions]?'' Explain why answering
this question is worthwhile.]

\section*{3. Proposed idea and novelty\wordbudget{110--140}}

[Describe exactly what you will change and what stays the same. Explain the inputs,
the main steps, and the expected output. Include a short concrete example or a small
equation only if it makes the idea clearer. Explain why the change might address the
gap. State how it differs from the originalp paper; do not present an existing paper
component as your new contribution. A modest, well-defined extension is sufficient.
Check related work before claiming broader novelty.]

\section*{4. Evaluation plan\wordbudget{120--140}}
[Name the task, dataset or test cases, model, and a feasible pilot size. Compare with
the original method and at least one simple alternative. Identify the main metric
and a concrete success criterion. State which resources and settings must be matched
for a fair comparison. Include one ablation that removes your proposed component.
Keep development and test examples separate. Explain how you will handle variability,
cost, failures, and an inconclusive result. For a theoretical idea, specify the claim,
assumptions, and proof or counterexample strategy instead.]

\section*{5. Expected outcomes and risks\wordbudget{50--70}}
[State your hypothesis, not invented results. Explain where the idea should help,
one plausible reason it could fail, and a fallback or diagnostic test. Say what a
negative result would teach us. Make clear which numerical values are goals rather
than observed improvements.]

\section*{6. Feasibility and deliverables\wordbudget{40--50}}
[Give a short schedule, required data or computing resources, and the smallest
version you could realistically test. Name the deliverables, such as a prototype,
comparison plot, proof sketch, or error analysis. Separate the essential pilot from
optional extensions.]

\vfill
\referencesheading
{\small\textbf{References.} [1] [Authors. \textit{Presented paper title}. Venue/year.
\href{https://example.com}{Paper link}.] [2] [Closest related work, if needed.]\par}
\end{document}
